本文研究地球同步轨道（GEO）卫星对 30 颗高轨目标的抵近观测机动规划问题。从二体轨道相对运动规律与光学观测几何条件出发，本文将距离球壳与双圆锥约束统一抽象为时变可观测管道，围绕先最大化观测目标数、再最小化累计速度增量的两级目标，构建了从单星单标、单星多标到双星协同逐层递进的机动规划框架，核心为\textbf{基于 Clohessy--Wiltshire（CW）动态规划的时空联合序列凸优化模型}。

\textbf{针对问题一}，核心思路是利用轨道几何特征缩小搜索范围，再由可行初值出发做局部精修。同步轨道目标的轨道面差异较大，直接大角度变轨代价高昂，因此优先利用两星轨道面交线与背阳侧光照条件的交集锁定低代价候选观测窗口，沿管道截面角平分线构造满足几何约束的两脉冲初值，再通过序列二次规划（SLSQP）对机动时刻、脉冲矢量与观测时刻联合求解。最终得到 48~h 内观测目标 B1 的最优两脉冲方案，累计速度增量为 \textbf{5.2128~m/s}，两次机动间隔 13.56~h。

\textbf{针对问题二}，30 天内单星 R1 巡游 30 颗目标的难点在于离散访问顺序与连续轨道机动的强耦合，为此采用\textbf{离散排序与连续精修分层协同}的思想。离散层利用三维 CW 相对运动近似刻画转移代价，通过物理状态动态规划统筹访问顺序与交会日期，自然利用目标分布规律减少折返并识别 B19 与 B20 的同次伴随飞越；连续层采用切平面凸化处理非凸距离下界，将脉冲矢量、机动时刻与观测起始时刻纳入序列凸优化同步精修，并对归零的冗余脉冲重新分配位置以进一步降低燃料消耗。最终实现 \textbf{30 颗目标全覆盖}，累计速度增量为 \textbf{133.03~m/s}，全部机动与观测约束均严格满足。

\textbf{针对问题三}，引入新星 R2 后规划的关键在于消除单星往返折返并充分发挥初始轨道设计的自由度。基于同步带目标相对经度分布，将西侧 10 颗与东侧 20 颗目标分别交由 R1 与 R2 双向单程推进，消除折返路程；在此基础上，将 R2 的初始六维轨道状态直接嵌入时空联合序列凸优化模型，实现\textbf{初始轨道与双星全程机动的同步设计}。最终两星合计实现 \textbf{30 颗目标全覆盖}，其中 R1 消耗 \textbf{11.94~m/s}，R2 消耗 \textbf{14.12~m/s}，R2 初始半长轴较同步轨道低 5.895~km、偏心率为 0.000950，两星累计速度增量合计 \textbf{26.07~m/s}。

\textbf{模型检验与灵敏度分析表明}，三问\textbf{全部 61 次观测}均通过高密度时间网格校验与 ATK 4.2.0 单星机动规划（MCS）连续数值积分复核，解析轨迹与数值积分最大位置差不超过 0.152~m，软件原生覆盖时段报告显示所有目标连续覆盖时长均超过 60~s；全程轨道面倾转角极小，印证了准共面分析的合理性；灵敏度分析进一步表明，问题三中\textbf{半长轴偏离上限是主导约束}，所得方案在安全裕量与轨道圆度变化下均保持良好稳健性。

\vspace{4pt}
\noindent{\heiti 关键词：}高轨卫星抵近观测；Clohessy--Wiltshire 相对运动；物理状态动态规划；序列凸优化；双星协同巡游
